The directory at \url{https://aibenchmarkindex.com} is a list, not a leaderboard. Each package prints an index entry. The fields we rely on are the id, name, tagline, tags, category, a short statement of what the benchmark measures, the npm name, the repository, a docs URL, the version, the task count, the data sources (name, kind, licence, URL), the grader string, the status, saturation, created and updated dates, and a share-card URL. \texttt{benchmark-core} exports that shape as \texttt{IndexEntry}. Table~\ref{tab:graders} is those strings, not a paraphrase.

The permalink contract is \texttt{/benchmarks/<package>}. The pelican docs URL in the entry is \url{https://aibenchmarkindex.com/benchmarks/benchmark-pelican}, and the other packages follow the same pattern. A directory page should keep that path when the package name does not change.

Discovery is two files in the package, both generated for the MCP registry and for agents that read a short brief. \texttt{server.json} names the stdio server. \texttt{llms.txt} is the short description. Neither file is a score. \texttt{npx <package> mcp} starts the server the registry entry points at.

Status in this release is \texttt{\suiteStatus{}} for every package. That is a warning. The tasks and graders are real enough to run, and they are not a stable evaluation standard. A breaking change to a grader should be a new version, because an old report would otherwise become incomparable. The index entry carries the version for that reason.
